\section*{Chapitre \ref{ch:pl:capturing-and-storing-tick-data} --- Capture et stockage des données tick}

\begin{solution}{exo:pl:capturing-and-storing-tick-data:1}
$15\,773\,858 / 258\,535 = 61.0$ octets par message. L'encapsulation représente $12 + 20 + 2 = 34$ octets par paquet, et avec un message par paquet
soit $34/61.0 = 55.7~\%$ de la capture~: plus de la moitié des octets sont des enveloppes.
\end{solution}

\begin{solution}{exo:pl:capturing-and-storing-tick-data:2}
$\mathrm{CR} = 56 / 15.24 = 3.67$.
\end{solution}

\begin{solution}{exo:pl:capturing-and-storing-tick-data:3}
$811\,160 / 56 = 14\,485$ enregistrements. Avec un \omterm{def:pl:capturing-and-storing-tick-data:partition}{partitionnement} par instrument, la requête lit $811\,160 / 14\,477\,960 = 5.6\%$ de la journée~;
sans lui, la journée entière (un fichier trié muni d'un index pourrait faire mieux, chapitre~3).
\end{solution}

\begin{solution}{exo:pl:capturing-and-storing-tick-data:4}
Seules les deux \omterm{def:pl:capturing-and-storing-tick-data:raw}{captures brutes} indiquent quelle ligne a perdu quel paquet et quand chaque copie est arrivée~: c'est la question d'une enquête sur le réseau ou sur la latence
(la ligne B était-elle plus lente ce matin~? la panne a-t-elle causé le trou~?). Toute question sur le marché lui-même (le carnet
d'ordres à 10~h~03, les transactions sur l'instrument~2) reçoit une réponse tout aussi bonne des enregistrements arbitrés, qui contiennent chaque message une seule fois.
\end{solution}

\begin{solution}{exo:pl:capturing-and-storing-tick-data:5}
$311\times10^9 \times 11.9 = 3.70\times10^{12}$ octets, soit 3.70~To par jour~; $\times 252 = 0.93$~PB par an.
\end{solution}

\begin{solution}{exo:pl:capturing-and-storing-tick-data:6}
Cinq ans représentent 4.67~Po, soit $4.67\times10^6$~GB~; à 0.023~USD le GB-mois, cela fait environ 1\,288\,000~USD par an, soit 6.6 fois le coût de la
politique à niveaux de stockage, de 194\,000~USD. Presque toute l'économie provient du transfert des quatre années les plus anciennes vers l'archive.
\end{solution}

\begin{solution}{exo:pl:capturing-and-storing-tick-data:7}
Le taux baisse légèrement, de 4.70 à 4.67. Trier par instrument réduit les sauts de la colonne des prix (plus de sauts entre les niveaux de prix
des deux instruments), mais il casse les colonnes régulières dans l'ordre d'arrivée~: les numéros de séquence n'augmentent plus
exactement d'une unité, les heures de réception et d'échange reculent à la frontière et progressent par des pas plus grands et irréguliers. Sur cette séance, les
deux effets se compensent presque. La leçon du chapitre est~: mesurer la disposition sur les données elles-mêmes.
\end{solution}

\begin{solution}{exo:pl:capturing-and-storing-tick-data:8}
Une capture doit suivre le débit de crête du flux, et non sa moyenne, et elle écrit chaque octet une seule fois. Sur le portable, zlib au niveau~9
compresse à 3.0~Mo/s, contre 34.5 au niveau~6 et 220 pour zstd au niveau~1 dans un fichier en colonnes, pour un \omterm{def:pl:capturing-and-storing-tick-data:ratio}{taux de compression} supérieur de 2~\% à celui du niveau~6~;
au pic d'un flux complet, il prendrait du retard en quelques secondes et perdrait des paquets, et les paquets par appel sont trop petits pour bien se compresser
de toute façon. Compresser la capture vite, par blocs, et recompresser la nuit pour gagner en taux.
\end{solution}

\begin{solution}{pb:pl:capturing-and-storing-tick-data:1}
\begin{enumerate}
\item 503\,024 paquets~; 258\,535 messages.
\item 255\,899~: chaque message est envoyé sur les deux lignes, si bien que presque chaque paquet a une seconde copie sur l'autre ligne~; l'écart
par rapport à des moitiés exactes tient aux paquets que chaque ligne a perdus et aux paquets dupliqués.
\item Huit messages dans quatre trous. Un feed handler en production les demanderait au service de retransmission ou reconstruirait le carnet à partir
du canal d'instantanés (snapshot).
\item L'arbitrage et la normalisation éliminent les doublons, la ligne qui a livré la première, les frontières des paquets, les
heures d'arrivée des copies suivantes et tout champ sans place dans l'enregistrement.
\item L'indicateur «~imprimable~» du message d'exécution avec prix (ou le symbole et le numéro de suivi de l'ordre d'ajout).
\item 61.0, 56.0 et 11.9 octets par message.
\item 3.32 et 4.70.
\item Dans l'ordre d'arrivée, il augmente exactement d'une unité à chaque enregistrement, si bien que ses deltas forment une constante qui se compresse presque
à rien.
\item La \omterm{def:pl:capturing-and-storing-tick-data:raw}{capture brute}, compressée rapidement, est gardée aussi longtemps que les enquêtes l'exigent (des semaines à des mois)~; les enregistrements en colonnes,
à encodage delta, sont gardés pour la durée de rétention (cinq ans).
\item L'encapsulation, la taille des messages et la redondance sont des propriétés du flux~; le simulateur envoie un message par paquet.
\item 37.95~To par jour.
\item 47.8~Po.
\item 3.70~To par jour, 0.93~Po par an, 4.67~Po sur cinq ans.
\item Environ 1\,989\,000~USD par an pour la \omterm{def:pl:capturing-and-storing-tick-data:raw}{capture brute} des deux lignes et 194\,000~USD pour les colonnes compressées.
\item Il s'agit de la capacité que prévoit l'opérateur du flux, un plafond et non une journée moyenne.
\item \textbf{Résultat nommé.} À la capacité prévue du flux d'options, la capture complète représente 9.56~Po par an sous forme de paquets bruts sur les deux lignes
(1.99~million~USD par an pour conserver cinq ans sur trois niveaux) et 0.93~Po par an sous forme de colonnes à encodage delta (0.19~million~USD)~: un
pétaoctet par an seulement dans la disposition la plus légère.
\item Pour les deux, dans des proportions comparables~: la disposition divise les octets par une dizaine par rapport à la \omterm{def:pl:capturing-and-storing-tick-data:raw}{capture brute} des deux lignes, le découpage par niveaux divise le
coût des mêmes octets par environ six par rapport au stockage standard.
\item Elle lit 5.6~\% de la journée au lieu de la totalité~; avec des millions de séries, un \omterm{def:pl:capturing-and-storing-tick-data:partition}{partitionnement} par série donne des millions de minuscules
fichiers dont l'ouverture coûte plus que la lecture, si bien que la clé doit être plus grossière (le sous-jacent, la date).
\item La compression rapide au moment de la capture, où le pic du flux fixe le rythme~; la compression lente, la nuit, quand la journée est réécrite une fois
pour des années de lecture.
\item La \omterm{def:pl:capturing-and-storing-tick-data:raw}{capture brute}, brièvement conservée pour les enquêtes, et les événements arbitrés pendant des années, en colonnes compressées et partitionnées, dans le
\omterm{def:pl:capturing-and-storing-tick-data:tier}{niveau de stockage} le moins cher que ses lecteurs peuvent tolérer.
\end{enumerate}
\end{solution}

\begin{solution}{iq:pl:capturing-and-storing-tick-data:1}
Les données brutes sont les octets tels que reçus, avec leurs heures d'arrivée et leurs doublons~; les données normalisées sont un événement par message dans le
format de la firme. Le normaliseur peut comporter des bogues ou manquer d'un champ, et seule la \omterm{def:pl:capturing-and-storing-tick-data:raw}{capture brute} permet de le relancer~; les données normalisées sont ce que
tout le monde lit.
\iqlookfor{Les données normalisées dépendent des données brutes, et non l'inverse.}
\end{solution}

\begin{solution}{iq:pl:capturing-and-storing-tick-data:2}
Fusionner par numéro de séquence~: la première copie l'emporte, les copies suivantes sont écartées~; laisser un trou ouvert un court instant, au cas où l'autre ligne
le comble~; puis demander au service de retransmission, ou reconstruire le carnet à partir d'un instantané (snapshot), et, dans une capture de recherche, consigner l'étendue du trou.
\iqlookfor{Numéros de séquence, doublons, délai d'attente, et reprise ou consignation du trou.}
\end{solution}

\begin{solution}{iq:pl:capturing-and-storing-tick-data:3}
Messages par jour (d'après les statistiques publiées par le flux et sa capacité prévue) multipliés par les octets par message dans chaque disposition (mesurés
sur un échantillon capturé, et non supposés), divisés par le \omterm{def:pl:capturing-and-storing-tick-data:ratio}{taux de compression} (mesuré sur le même échantillon), multipliés par le nombre de copies et par les jours de
négociation et les années~; puis en calculer le coût par niveau et par durée de rétention. Mesurer d'abord les octets par message et le \omterm{def:pl:capturing-and-storing-tick-data:ratio}{taux de compression}.
\iqlookfor{La mesure plutôt que l'hypothèse~; copies et niveaux comptés.}
\end{solution}

\begin{solution}{iq:pl:capturing-and-storing-tick-data:4}
D'après les filtres qu'utilisent les requêtes~: presque toujours la date, puis l'instrument ou le sous-jacent. Une clé trop fine crée de nombreux petits fichiers,
chacun ayant un coût fixe de listage et d'ouverture, si bien que les requêtes qui couvrent beaucoup de clés deviennent lentes et que le système de fichiers est mis à rude épreuve.
\iqlookfor{Élagage face au surcoût par fichier.}
\end{solution}

\begin{solution}{iq:pl:capturing-and-storing-tick-data:5}
Des messages consécutifs partagent l'essentiel de leur contenu~: les mêmes instruments, des prix voisins, des heures qui progressent par petits pas,
des numéros de séquence qui augmentent d'une unité. Stockés en colonnes, avec des différences pour les champs monotones et lentement variables, ces valeurs deviennent
de longues suites de petites valeurs~; les numéros de séquence et les horodatages se compressent le mieux, les références d'ordre le moins bien.
\iqlookfor{Colonnes et deltas~; quels champs sont réguliers.}
\end{solution}

\begin{solution}{iq:pl:capturing-and-storing-tick-data:6}
Capture filaire des deux lignes sur le site de colocation, conservée des semaines, avec une compression rapide~; événements normalisés et arbitrés, écrits au fil de la journée
dans l'ordre d'arrivée~; réécriture nocturne en colonnes partitionnées par date et sous-jacent, compressées au meilleur taux, conservées cinq ans ou
davantage~; niveau chaud pour le mois récent, accès peu fréquent pour l'année, archive au-delà. Dimensionner à partir de la capacité prévue du flux et
des octets par message mesurés~: des pétaoctets sur la durée de rétention, un coût de l'ordre de centaines de milliers à quelques millions par an selon
la disposition et les niveaux de stockage.
\iqlookfor{Point de capture, disposition, \omterm{def:pl:capturing-and-storing-tick-data:partition}{partitionnement}, niveaux et ordre de grandeur du coût avec ses hypothèses.}
\end{solution}
